
\section{Mixing from an arbitrary polymer configuration}\label{sec:full-polymer-tv}
\begin{theorem}\label{thm:full-polymer-tv-upper-new}
For every fixed $b>0$, the polymer chain with $\beta=b/n$ satisfies
\[
 t^{\mathrm{polymer}}_{\mathrm{mix},n}(\varepsilon)
 \leq C_b\sqrt n\left\{(\log n)^2+\log(1/\varepsilon)\right\},
 \qquad 0<\varepsilon<1,
\]
for all sufficiently large $n$. Together with endpoint cutoff this gives,
for each fixed $\varepsilon\in(0,1)$,
\[
 \frac{\sqrt n\log n}{8\sqrt b}-O_{b,\varepsilon}(\sqrt n)
 \leq t^{\mathrm{polymer}}_{\mathrm{mix},n}(\varepsilon)
 \leq O_b\bigl(\sqrt n(\log n)^2\bigr).
\]
The extra logarithm in the upper bound is not removed by this argument;
worst-case polymer cutoff is not asserted.
\end{theorem}

\subsection{A degree-dependent sector bound}
Let $L_B$ be the labelled-cut comparison generator, and let
$\mathcal L_\ell$ be its degree-$\ell$ radial generator:
\begin{align*}
 (\mathcal L_\ell g)(k)
 ={}&\beta(n-k-\ell)(g(k+1)-g(k))\\
 &+\frac{(k-1)(k-\ell)}n(g(k-1)-g(k))
 -\left(\ell\beta+\frac{\ell(k-1)}n\right)g(k),
 \qquad \ell\leq k\leq n-\ell.
\end{align*}
Assume $2\leq A:=n\sqrt\beta\leq n/4$. For $\ell\geq1$ use
\[
 h_\ell(k)=\frac1{(k+A)_\ell},\qquad
 (x)_\ell=x(x+1)\cdots(x+\ell-1).
\]
Its neighboring ratios are
\[
 \frac{h_\ell(k+1)}{h_\ell(k)}=\frac{k+A}{k+A+\ell},\qquad
 \frac{h_\ell(k-1)}{h_\ell(k)}=\frac{k+A+\ell-1}{k+A-1}.
\]
Direct substitution gives the exact identity
\begin{equation}\label{eq:degree-certificate-new}
 \frac{-(\mathcal L_\ell h_\ell)(k)}{h_\ell(k)}
 =\ell\left[
 \frac{\beta(n-k-\ell)}{k+A+\ell}+\beta
 +\frac{(k-1)(A+\ell-1)}{n(k+A-1)}\right].
\end{equation}
All bracketed terms are nonnegative. If $k\leq A$, then $\ell\leq k$
gives $n-k-\ell\geq n/2$ and $k+A+\ell\leq3A$; the first term is
at least $\sqrt\beta/6$. If $k\geq A$, then
\[
 \frac{k-1}{k+A-1}\geq\frac13,\qquad A+\ell-1\geq A,
\]
so the last term is at least $\sqrt\beta/3$.
The finite positive-vector comparison therefore proves that every
eigenvalue in degree $\ell$ is at least
\begin{equation}\label{eq:degree-linear-lower-new}
 \lambda_\ell\geq\frac{\ell\sqrt\beta}{6}.
\end{equation}
Independently, the killing term gives
\[
 \lambda_\ell\geq\ell\beta+\frac{\ell(\ell-1)}n.
\]
In particular, at fixed $b>0$, $\beta=b/n$, and $A=\sqrt{bn}$,
\[
 \lambda_\ell\geq \frac{b\ell}{6A}
 \qquad(1\leq\ell\leq\lfloor n/2\rfloor)
\]
for every sufficiently large $n$.


\subsection{A polynomial bound on the count density}
Fix $b>0$, put $A=\sqrt{bn}$, and let $p$ be the stationary count law.
Let $r$ be a mode of $p$ and define
\[
 V(k)=\left(\frac{p(r)}{p(k)}\right)^{3/4}\geq1,
 \qquad W(k)=(\sqrt{\lambda_k}-\sqrt{\mu_k})^2.
\]
We prove the drift inequality
\begin{equation}\label{eq:count-log-lyapunov-new}
 \frac{(QV)(k)}{V(k)}
 \leq\frac{C_b-c_b\log V(k)}{A},\qquad 1\leq k\leq n,
\end{equation}
for all sufficiently large $n$.

By detailed balance, the exact expression is
\[
 \frac{QV(k)}{V(k)}
 =\lambda_k^{1/4}\mu_{k+1}^{3/4}
  +\mu_k^{1/4}\lambda_{k-1}^{3/4}-\lambda_k-\mu_k.
\]
At an endpoint the term with zero outgoing rate is zero; thus the displayed
formula remains valid if one formally sets $\lambda_0=b$ and
$\mu_{n+1}=n+1$. These are algebraic conventions only.
Write $\lambda=\lambda_k$ and $\mu=\mu_k$. With
$u=\lambda^{1/4}$ and $v=\mu^{1/4}$,
\begin{align*}
 &\lambda+\mu-\lambda^{1/4}\mu^{3/4}-\mu^{1/4}\lambda^{3/4}\\
 &\hspace{1cm}=(u-v)^2(u^2+uv+v^2)
 \geq\frac34(u-v)^2(u+v)^2
 =\frac34W(k).
\end{align*}
The two coefficient-variation errors are
\[
 E_1=\lambda_k^{1/4}(\mu_{k+1}^{3/4}-\mu_k^{3/4}),\qquad
 E_2=\mu_k^{1/4}(\lambda_{k-1}^{3/4}-\lambda_k^{3/4}).
\]
Since $\mu_{k+1}-\mu_k=2k/n$, concavity gives, for $k\geq2$,
\[
 E_1\leq\frac34\lambda_k^{1/4}\mu_k^{-1/4}\frac{2k}{n}
       \leq C_b n^{-3/4}\sqrt k.
\]
The same bound holds at $k=1$ by direct substitution, since
$E_1=\lambda_1^{1/4}(2/n)^{3/4}$.
For $k\leq2A$, it follows that $E_1\leq C_b/A$.
For $k>2A$, one has $\mu_k\geq3b\geq3\lambda_k$ and hence
$W(k)\geq c_bk^2/n$ (for large $n$). Thus
\[
 \frac{E_1}{W(k)}\leq C_bn^{1/4}k^{-3/2}\leq\frac{C_b}{A}.
\]
For sufficiently large $n$, this is at most $1/4$.
The inequality $(x+y)^{3/4}-x^{3/4}\leq y^{3/4}$ also gives, uniformly
including $k=n$,
\[
 E_2\leq\mu_k^{1/4}(b/n)^{3/4}
 \leq\frac{b^{3/4}\sqrt k}{n}\leq\frac{C_b}{A}.
\]
Consequently
\begin{equation}\label{eq:count-drift-square-new}
 \frac{QV(k)}{V(k)}\leq-\frac12W(k)+\frac{C_b}{A}.
\end{equation}

It remains to check
\begin{equation}\label{eq:stationary-weight-vs-drift-new}
 \log V(k)\leq C_b\{1+A W(k)\}.
\end{equation}
On $A/2\leq k\leq2A$, the elementary stationary ratio estimates proved
earlier give
\[
 \log\frac{p(r)}{p(k)}\leq C_b\left(1+\frac{(k-A)^2}{A}\right).
\]
The exact drift factorization around its zero $a=A+O_b(1)$ and
$\lambda_k+\mu_k\leq C_b$ give
$W(k)\geq c_b(k-a)^2/A^2$, proving
\eqref{eq:stationary-weight-vs-drift-new} on this interval.
If $k<A/2$, then $W(k)\geq c_b>0$, while
\[
 \log\frac{p(r)}{p(k)}\leq\log\frac{p(r)}{p(1)}
 \leq\log\frac{A^{2(r-1)}}{(r-1)!\,r!}\leq C_bA,
\]
where $r=A+O_b(1)$ and the last bound follows from
$j!\geq(j/e)^j$.
If $k>2A$, sum the consecutive ratios to obtain
\[
 \log\frac{p(r)}{p(k)}
 =\sum_{j=r}^{k-1}
 \left[\log\frac jA+\log\frac{j+1}A-\log(1-j/n)\right]
 \leq\frac{C_bk^2}{A}.
\]
Here $\log x\leq x$ controls the first two terms. For the last,
$-\log(1-j/n)\leq2j/n$ when $k\leq n/2$; when $k>n/2$, its sum is
at most
$\sum_{j=1}^{n-1}-\log(1-j/n)\leq Cn\leq Ck^2/n$.
Since $W(k)\geq c_bk^2/n$ and $A^2=bn$, this proves
\eqref{eq:stationary-weight-vs-drift-new} in the remaining case.
Combining it with \eqref{eq:count-drift-square-new} proves
\eqref{eq:count-log-lyapunov-new}.

Now let $y(t)=\mathbb E_iV(K_t)$ for any initial count $i$.
The finite-state evolution equation and Jensen's inequality for
$x\log x$ imply
\[
 y'(t)\leq\frac{C_b}{A}y(t)-\frac{c_b}{A}y(t)\log y(t).
\]
Thus $z(t)=\log y(t)$ satisfies
\[
 z(t)\leq\frac{C_b}{c_b}+e^{-c_bt/A}z(0).
\]
Uniformly in $i$, \eqref{eq:stationary-weight-vs-drift-new} and
$W(i)\leq\lambda_i+\mu_i\leq b+n$ give
$z(0)\leq C_bn^2$. Consequently there exists $C_b<\infty$ such that
\begin{equation}\label{eq:count-V-moment-new}
 \sup_i\sup_{t\geq t_0}\mathbb E_iV(K_t)\leq C_b,
 \qquad t_0=C_b A\log n.
\end{equation}

Let $q_t$ be any count distribution at time $t\geq t_0$ started from a
point mass (or any mixture). Equation \eqref{eq:count-V-moment-new}
gives $q_t(k)\leq C_b/V(k)$. The stationary Gaussian and geometric-tail
estimates give
\[
 p(r)\asymp_b A^{-1/2},\qquad \sum_k\sqrt{p(k)}\leq C_b A^{1/4}.
\]
It follows that
\begin{equation}\label{eq:count-L2-polynomial-new}
 M(t):=\sum_{k=1}^n\frac{q_t(k)^2}{p(k)}
 \leq C_bp(r)^{-3/2}\sum_k\sqrt{p(k)}\leq C_bA,
 \qquad t\geq t_0.
\end{equation}
This is a polynomial bound, not a claim that count $L^2$ mixing is complete.


\subsection{The polymer mixing bound}


Let $p_{\mathrm{cut}}$ be the map from nonempty labelled cut sets to their
unordered segment lengths, and write $\pi_B$, $\pi_M$, and $p$ for the
stationary cut, polymer, and count laws, respectively. Start the actual
polymer chain from any state $x$. Let $f_t$ be its density with respect to
$\pi_M$, and lift it by
\[
 F_t(S)=f_t(p_{\mathrm{cut}}(S)).
\]
The lift is nonnegative. Its average on the layer $|S|=k$ is
\[
 \overline F_t(k)=\frac{q_t(k)}{p(k)},
\]
where $q_t$ is the distribution of the actual polymer count. Strong
lumpability of the count means that $q_t$ evolves under $Q$ from $K(x)$.
The preceding Lyapunov calculation therefore gives, with $A=\sqrt{bn}$,
\begin{equation}\label{eq:polymer-count-warm-new}
 M(t):=\sum_kp(k)\overline F_t(k)^2\leq M_*:=C_bA,
 \qquad t\geq t_0:=C_bA\log n.
\end{equation}
Choose the constant so that $M_*\geq1$.

Let $P_\ell$ denote the orthogonal projection onto the complete
degree-$\ell$ sector of the labelled-cut comparison space, and set
\[
 d_\ell=\dim H_\ell=\binom n\ell-\binom n{\ell-1}.
\]
On a layer of size $N_k=\binom nk$, the projector onto degree $\ell$
has constant diagonal $d_\ell/N_k$ whenever that degree is present.
Indeed, the projector commutes with every permutation of the labels,
and its trace is $d_\ell$. Positive semidefiniteness then gives
\[
 |P_\ell(S,T)|\leq\sqrt{P_\ell(S,S)P_\ell(T,T)}=d_\ell/N_k.
\]
For every nonnegative function $F$, it follows, using the ordinary matrix
kernel of this layerwise projector, that
\begin{align*}
 \|P_\ell F\|_{L^2(\pi_B)}^2
 &=\sum_k\frac{p(k)}{N_k}
       \sum_{|S|=|T|=k}F(S)P_\ell(S,T)F(T)\\
 &\leq d_\ell\sum_kp(k)\overline F(k)^2.
\end{align*}
Absent degrees contribute zero. Telescoping the dimensions gives
\begin{equation}\label{eq:positive-low-degree-new}
 \sum_{\ell=0}^{d}\|P_\ell F_t\|_2^2
 \leq\binom nd M(t),\qquad 0\leq d\leq\lfloor n/2\rfloor.
\end{equation}

Put $R(t)=\|F_t\|_{L^2(\pi_B)}^2=\|f_t\|_{L^2(\pi_M)}^2$ and
$g=\sqrt\beta/6=b/(6A)$. The stationary Dirichlet-form identity and
the proved sector bound $\lambda_\ell\geq g\ell$ imply
\begin{align*}
 -R'(t)
 &=2\mathcal E_M(f_t,f_t)=2\mathcal E_B(F_t,F_t)\\
 &\geq2g(d+1)\sum_{\ell>d}\|P_\ell F_t\|_2^2\\
 &\geq2g(d+1)\left\{R(t)-M_*\binom nd\right\},
 \qquad t\geq t_0.
\end{align*}
This is a statement about the actual polymer evolution; no invariance of
the polymer lift under the cut semigroup, or under individual sector
projections, is assumed. The sector decomposition is used only to bound
the instantaneous stationary energy.

Positivity also bounds the initial size of this differential inequality:
on each layer, $\sum_SF_t(S)^2\leq(\sum_SF_t(S))^2$, so
\[
 R(t)\leq\sum_kp(k)\binom nk\overline F_t(k)^2
       \leq2^nM_*,\qquad t\geq t_0.
\]
For $n$ large, whenever $R(t)>16M_*n^2$, put
\[
 s(t)=\log\frac{R(t)}{M_*},\qquad
 d=\left\lfloor\frac{s(t)}{2\log n}\right\rfloor.
\]
Then $1\leq d<\lfloor n/2\rfloor$, because
$s(t)\leq n\log2$. Moreover,
\[
 M_*\binom nd\leq M_*n^d\leq M_*e^{s(t)/2}\leq R(t)/4.
\]
The energy inequality consequently gives
\[
 s'(t)=\frac{R'(t)}{R(t)}
 \leq-g(d+1)\leq-\frac{g}{2\log n}s(t).
\]
Since $s(t_0)\leq n\log2$, it takes at most
$2g^{-1}(\log n)^2$ additional time to reach
$R(t)\leq16M_*n^2$. Once this threshold is reached it remains satisfied,
because $R(t)$ is nonincreasing.

Finally use the already established polymer gap
$\gamma_M\geq c_b/A$. If $t_1$ is the preceding threshold time, then
\[
 R(t_1+u)-1\leq e^{-2\gamma_Mu}\{R(t_1)-1\},\qquad
 \|P^M_{t_1+u}(x,\cdot)-\pi_M\|_{\rm TV}
 \leq\frac12\sqrt{R(t_1+u)-1}.
\]
As $R(t_1)\leq16C_bAn^2$, an additional
$C_bA\{\log n+\log(1/\varepsilon)\}$ time suffices. This proves
Theorem~\ref{thm:full-polymer-tv-upper-new}, uniformly over the initial
polymer state $x$.
